\documentclass[czech, oth, kiv, he, pdf, viewonly]{fasthesis}

\usepackage{float}

\title{Skladový systém}
\worktypespec{Dokumentace semestrální práce}
\author{Petr}{Schöpp}{}{}

\nobastardtitle
\nocopyrightnotice

\begin{document}

\frontpages[notm]
\tableofcontents

\chapter{Úvod}

Tato semestrální práce se zabývá návrhem a implementací
skladového systému. Aplikace je určena pro evidenci skladových
položek a skladových pohybů, především příjmů a výdejů.

Motivací pro vytvoření aplikace byla potřeba jednoduchého nástroje,
který může skladník používat přímo ve skladu. Proto je hlavní část
systému řešena jako mobilní aplikace pro Android.

Cílem práce bylo vytvořit aplikaci, která umožní přehledně zobrazit
aktuální stav skladu, vyhledávat položky, vytvářet skladové doklady
a sledovat historii provedených pohybů.

Systém je rozdělen na mobilní aplikaci, serverovou část a databázi.
Mobilní aplikace zajišťuje uživatelské rozhraní a komunikaci se
serverem. Serverová část je vytvořena v prostředí Node.js a poskytuje
REST API. Data jsou ukládána do databáze PostgreSQL.

Výsledná aplikace podporuje práci se skladovými položkami, vytváření
příjmových a výdejových dokladů, skenování čárových a QR kódů,
ruční vyhledávání produktů, evidenci historie a ukládání fotografií
produktů.

\chapter{Popis aplikace}

\section{Účel aplikace}

Účelem aplikace je zjednodušit evidenci skladových položek a
skladových pohybů. Aplikace je navržena tak, aby ji mohl skladník
používat přímo ve skladu na běžném mobilním telefonu s operačním
systémem Android.

Důležitou motivací bylo také snížení nákladů. Specializovaná
skladová čtecí zařízení mohou být pro menší provozy zbytečně drahá.
Kromě samotného zařízení může být u některých komerčních řešení
nutné platit také licenci nebo pravidelné poplatky za používání
aplikace. Běžný mobilní telefon s Androidem je oproti tomu dostupnější
a často již obsahuje vše potřebné pro základní práci ve skladu,
například fotoaparát pro skenování kódů.

Aplikace proto využívá mobilní telefon jako jednoduché skladové
zařízení. Uživatel může zobrazit aktuální stav skladu, vyhledat
položku podle názvu nebo kódu, zobrazit detail položky a vytvořit
příjmový nebo výdejový doklad.

Při příjmu je možné založit novou položku, pokud se ve skladu ještě
nenachází. Při výdeji aplikace kontroluje, aby nebylo možné vydat
větší množství, než je aktuálně dostupné na skladě.

\section{Hlavní části systému}

Systém je rozdělen na tři hlavní části: mobilní aplikaci, server a
databázi.

Mobilní aplikace zajišťuje uživatelské rozhraní, práci s formuláři,
skenování kódů a komunikaci se serverem. Server zpracovává požadavky
mobilní aplikace, provádí validaci dat a ukládá změny do databáze.

Databáze uchovává informace o uživatelích, skladech, položkách,
kódech položek, skladových množstvích a skladových dokladech.

\section{Použité technologie}

Aplikace je rozdělena na tři hlavní části: mobilní aplikaci, server
a databázi. Každá část používá technologie odpovídající své úloze.

\subsection{Mobilní aplikace}

Mobilní aplikace je vytvořena v jazyce Kotlin pro operační systém
Android. Uživatelské rozhraní je vytvořeno pomocí Jetpack Compose.
Obrazovky jsou skládány z menších Composable funkcí, což usnadňuje
jejich opakované použití a pozdější úpravy.

Pro opakované části rozhraní byly vytvořeny vlastní komponenty.
Mezi ně patří například společná karta \texttt{AppCard}, seznamová
karta \texttt{AppListCard}, spodní akční lišta, obrazovka se skenerem
a komponenta pro barevné označení typu skladového pohybu. Tyto
komponenty sjednocují vzhled aplikace a omezují opakování kódu.

Aplikace využívá běžné prvky uživatelského rozhraní, například textová
pole, tlačítka, dialogy, rozbalovací menu, modální spodní panel a
seznamy. Důležité akce jsou umístěny hlavně do spodní části obrazovky,
aby byly na telefonu dobře dostupné.

Stav obrazovek je spravován pomocí ViewModelů. Stav je uchováván
pomocí \texttt{StateFlow} a v uživatelském rozhraní je pozorován pomocí
funkce \texttt{collectAsState}. Při změně dat se tak příslušná část
obrazovky automaticky překreslí.

Komunikace se serverem je řešena knihovnou Retrofit. Retrofit převádí
volání funkcí v mobilní aplikaci na HTTP požadavky na REST API serveru.
Data přenášená mezi serverem a aplikací jsou ve formátu JSON. Pro
mapování názvů polí z JSON odpovědí na Kotlin modely jsou použity
anotace \texttt{SerializedName}.

Pro načítání obrázků produktů je použita knihovna Coil. Obrázky jsou
uloženy na serveru a mobilní aplikace je načítá podle URL adresy.

Kamera je v aplikaci využita dvěma způsoby. Prvním je pořízení
fotografie produktu. Druhým je skenování čárových a QR kódů. Pro práci
s kamerou při skenování je použita knihovna CameraX. Náhled kamery je
zobrazen pomocí komponenty \texttt{PreviewView}, která je do Compose
rozhraní vložena přes \texttt{AndroidView}.

Rozpoznávání kódů je řešeno pomocí ML Kit Barcode Scanning. Obraz
z kamery je zpracováván pomocí \texttt{ImageAnalysis}. Jednotlivé
snímky jsou převedeny na \texttt{InputImage} a následně předány
scanneru z ML Kit. Pokud je rozpoznán čárový nebo QR kód, aplikace jej
použije pro vyhledání položky nebo pro přidání položky do příjmového
nebo výdejového dokladu.

Součástí skeneru je také vlastní vizuální překryv. Ten vykresluje
oblast rozpoznaného kódu a informuje uživatele, zda byl kód nalezen.
Pokud je v obrazu více kódů, aplikace vybírá kód nejblíže středu
obrazu, aby bylo skenování pro uživatele předvídatelnější.

\subsection{Serverová část}

Serverová část je vytvořena v prostředí Node.js s frameworkem Express.
Server poskytuje REST API, které využívá mobilní aplikace.

Server je rozdělen do samostatných rout podle oblastí systému.
Samostatně je řešeno přihlášení, položky, doklady, pohyby, uživatelé
a sklady. Díky tomu není serverová logika soustředěna v jednom velkém
souboru.

Pro přístup k databázi je použita knihovna \texttt{pg}. Server provádí
SQL dotazy nad databází PostgreSQL a vrací mobilní aplikaci data ve
formátu JSON.

Pro nahrávání fotografií produktů je použita knihovna Multer.
Fotografie se uloží do adresáře na serveru a v databázi se uloží název
souboru. Mobilní aplikace potom fotografii načítá přes veřejnou URL
adresu serveru.

Konfigurace serveru je načítána ze souboru \texttt{.env} pomocí
knihovny \texttt{dotenv}. Knihovna \texttt{cors} umožňuje komunikaci
mobilní aplikace se serverem v lokální síti.

\subsection{Databáze}

Data jsou ukládána do databáze PostgreSQL. Databáze uchovává informace
o uživatelích, skladech, skladových položkách, kódech položek,
aktuálním množství položek na skladě a skladových dokladech.

Databázové schéma je definováno v souboru \texttt{schema.sql}.
Testovací data jsou vložena pomocí souboru \texttt{seed.sql}.

Databáze využívá primární a cizí klíče pro zachování vazeb mezi
tabulkami. Například položky dokladů jsou navázány na konkrétní
skladový doklad a konkrétní skladovou položku. Tabulka
\texttt{warehouse\_items} uchovává aktuální množství položky na
konkrétním skladu.

\section{Architektura aplikace}

Aplikace je navržena jako klient-server systém. Mobilní aplikace
nepracuje přímo s databází, ale komunikuje se serverem pomocí HTTP
požadavků. Server následně provádí databázové operace a vrací
mobilní aplikaci odpovědi ve formátu JSON.

Toto rozdělení umožňuje oddělit uživatelské rozhraní od databázové
a serverové logiky. Mobilní aplikace řeší především zobrazení dat,
práci s uživatelským vstupem a komunikaci se serverem. Serverová část
řeší validaci požadavků, ukládání dat a změny skladového množství.

\subsection{Mobilní část}

Mobilní část je rozdělena do několika vrstev. Datové modely popisují
objekty, se kterými aplikace pracuje, například položky, doklady,
uživatele nebo skladové pohyby. Síťová vrstva obsahuje API klienta
a definice volání serverových endpointů.

Stav jednotlivých obrazovek je uložen ve ViewModelech. ViewModely
obsahují logiku pro načítání dat, ukládání změn a zpracování chyb.
Uživatelské rozhraní je tvořeno Composable funkcemi, které pouze
zobrazují aktuální stav a předávají uživatelské akce zpět do
ViewModelů.

Toto rozdělení zjednodušuje údržbu aplikace. Pokud se změní způsob
načítání dat ze serveru, není nutné měnit samotné vykreslení
obrazovek. Stejně tak lze upravovat vzhled komponent bez zásahu do
serverové komunikace.

\subsection{Serverová část}

Serverová část je rozdělena podle hlavních oblastí systému. Samostatné
routy řeší přihlášení, položky, doklady, pohyby, uživatele a sklady.
Každá routa obsahuje endpointy související s danou oblastí.

Například práce s položkami je oddělena od práce se skladovými
doklady. Díky tomu je kód přehlednější a nové funkce lze přidávat do
konkrétní části serveru bez nutnosti upravovat celý server.

Na serveru jsou také vyčleněny společné pomocné funkce. Ty slouží
například pro čtení parametrů z požadavku nebo pro jednotné vytváření
HTTP chyb. Díky tomu se neopakuje stejná logika v každé routě.

\subsection{Databázová část}

Databáze obsahuje samostatné tabulky pro uživatele, sklady, položky,
kódy položek, množství položek na skladě a skladové doklady. Vazby
mezi tabulkami jsou řešeny pomocí cizích klíčů.

Aktuální množství položek na skladě je uloženo v tabulce
\texttt{warehouse\_items}. Skladové doklady jsou uloženy v tabulkách
\texttt{stock\_documents} a \texttt{stock\_document\_items}. Díky tomu
je možné zvlášť evidovat aktuální stav skladu a zároveň uchovávat
historii provedených příjmů a výdejů.

Při vytvoření nebo úpravě dokladu server aktualizuje odpovídající
skladové množství. Tato logika je ponechána na serveru, aby nemohlo
dojít k nekonzistentním změnám pouze na straně mobilní aplikace.

\chapter{Datový model}

Datový model je navržen tak, aby odděloval skladové položky, jejich
kódy, aktuální množství na skladě a historii skladových dokladů.
Díky tomu lze uchovávat nejen aktuální stav skladu, ale také zpětně
dohledat, jakými příjmy a výdeji byl tento stav vytvořen.

Databázové schéma je definováno v souboru \texttt{schema.sql}.
Základní testovací data jsou vložena pomocí souboru \texttt{seed.sql}.

\section{Přehled tabulek}

Následující tabulka shrnuje hlavní databázové tabulky a jejich význam.

\begin{table}[H]
\centering
\begin{tabular}{|l|p{9cm}|}
\hline
\textbf{Tabulka} & \textbf{Význam} \\
\hline
\texttt{roles} &
Uživatelské role, například skladník nebo administrátor. \\
\hline
\texttt{users} &
Uživatelé aplikace, jejich přihlašovací údaje, jméno, příjmení,
role a informace o aktivitě účtu. \\
\hline
\texttt{warehouses} &
Sklady, jejich názvy, kódy, umístění a informace o aktivitě. \\
\hline
\texttt{items} &
Skladové položky, jejich název, jednotka, poznámka a název souboru
s fotografií. \\
\hline
\texttt{item\_codes} &
Kódy skladových položek. Jedna položka může mít více různých kódů. \\
\hline
\texttt{code\_types} &
Typy kódů, například EAN, QR nebo interní kód. \\
\hline
\texttt{warehouse\_items} &
Aktuální množství konkrétní položky na konkrétním skladu,
včetně umístění a minimálního množství. \\
\hline
\texttt{movement\_types} &
Typy skladových pohybů a jejich směr, například příjem nebo výdej. \\
\hline
\texttt{stock\_documents} &
Hlavičky skladových dokladů, například číslo dokladu, sklad,
typ pohybu, stav a uživatelé. \\
\hline
\texttt{stock\_document\_items} &
Jednotlivé položky skladových dokladů, tedy položka, množství
a poznámka. \\
\hline
\end{tabular}
\caption{Přehled hlavních databázových tabulek}
\end{table}

\section{Uživatelé a sklady}

Uživatelé systému jsou uloženi v tabulce \texttt{users}. Každý
uživatel má uživatelské jméno, heslo, jméno, příjmení, roli a informaci
o tom, zda je aktivní.

Role uživatelů jsou uloženy v tabulce \texttt{roles}. V aktuální verzi
aplikace jsou připraveny role skladníka a administrátora.

Sklady jsou uloženy v tabulce \texttt{warehouses}. Každý sklad má
název, kód, volitelné umístění a informaci o aktivitě. Aktuálně
aplikace pracuje hlavně s hlavním skladem, ale datový model umožňuje
evidovat více skladů.

\section{Skladové položky}

Skladové položky jsou uloženy v tabulce \texttt{items}. Položka
obsahuje název, jednotku, volitelný název souboru s fotografií,
poznámku a informaci o tom, zda je aktivní.

Kódy položek jsou odděleny do tabulky \texttt{item\_codes}. Jedna
položka tak může mít více různých kódů. To umožňuje například evidovat
EAN kód, QR kód nebo interní kód položky.

Typy kódů jsou uloženy v tabulce \texttt{code\_types}. Díky tomu není
typ kódu uložen pouze jako text u každé položky, ale je spravován
samostatně.

Aktuální množství položky na konkrétním skladu je uloženo v tabulce
\texttt{warehouse\_items}. Tato tabulka propojuje sklad a položku.
Kromě aktuálního množství obsahuje také umístění položky ve skladu
a minimální množství.

\section{Skladové doklady}

Skladové doklady jsou uloženy v tabulce \texttt{stock\_documents}.
Doklad obsahuje číslo dokladu, sklad, typ pohybu, stav, poznámku,
autora vytvoření a uživatele, který doklad naposledy upravil.

Jednotlivé položky dokladu jsou uloženy v tabulce
\texttt{stock\_document\_items}. Každý záznam obsahuje odkaz na doklad,
odkaz na skladovou položku, množství a případnou poznámku.

Typy skladových pohybů jsou uloženy v tabulce
\texttt{movement\_types}. Každý typ pohybu má kód, název a směr
pohybu. Směr pohybu určuje, zda se množství na skladě zvyšuje, nebo
snižuje.

V aplikaci jsou nejdůležitější typy pohybu příjem a výdej. Příjem
zvyšuje množství položky na skladě, zatímco výdej ho snižuje.

\section{Aktuální stav a historie}

Aktuální stav skladu je uložen v tabulce \texttt{warehouse\_items}.
Historie změn je odvozena ze skladových dokladů a jejich položek.

Toto rozdělení bylo zvoleno proto, aby bylo možné rychle zobrazit
současný stav skladu a zároveň uchovávat historii provedených operací.
Při vytvoření nebo úpravě dokladu server aktualizuje aktuální množství
ve skladu a současně ponechává uložený doklad v historii.

Díky tomu lze v detailu položky zobrazit nejen aktuální množství, ale
také seznam pohybů, které se dané položky týkají.



\end{document}